\Article{18}

Methodus art. praec. valores absolutos functionum $T$, $U$ suppeditavit, conditionesque assignavit, sub quibus signa aequalia vel opposita illis tribuenda sint: sed signa ipsa hinc nondum determinantur.\quad Hoc pro eo casu, ubi statuitur $k=1$, sequenti modo supplebimus.

Statuamus $\rho=\cos\frac12\omega+i\sin\frac12\omega$, ita ut fiat $r=\rho\rho$, patetque, propter $\rho^n=-1$ aequationem $[8]$ ita exhiberi posse

\[
\begin{aligned}
W&=2(1+\rho^{n-2})(1+\rho^{-4})(1+\rho^{n-6})(1+\rho^{-8})\\&\qquad\ldots(1+\rho^{-n+4})(1+\rho^2)
\end{aligned}
\]

sive factoribus alio ordine dispositis

\[
\begin{aligned}
W&=2(1+\rho^2)(1+\rho^{-4})(1+\rho^6)(1+\rho^{-8})\\&\qquad\ldots(1+\rho^{-n+4})(1+\rho^{n-2})
\end{aligned}
\]

Iam fit

\[
\begin{aligned}
1+\rho^2&=2\rho\cos\frac12\omega\\
1+\rho^{-4}&=2\rho^{-2}\cos\omega\\
1+\rho^{-6}&=2\rho^3\cos\frac32\omega\\
1+\rho^{-8}&=2\rho^{-4}\cos2\omega,\;\text{etc.}
\end{aligned}
\]

usque ad

\[
\begin{aligned}
1+\rho^{-n+4}&=2\rho^{-\frac12n+2}\mathord{,}\cos(\frac14n-1)\omega\\
1+\rho^{n-2}&=2\rho^{\frac12n-1}\cos(\frac14n-\frac12)\omega
\end{aligned}
\]

Quamobrem habetur

\SourceLeaf{41}

\[
\begin{aligned}
W&=2^{\frac12n}\rho^{\frac14n}\cos\frac12\omega\cos\omega\,\frac32\omega\\&\qquad\dotsfour\cos(\frac14n-\frac12)\omega
\end{aligned}
\]

Cosinus in hoc productum ingredientes manifesto omnes positivi sunt, factor $\rho^{\frac14n}$ autem fit $=\cos45^0+i\sin45^0=(1+i)\surd\frac12$.\quad Hinc colligimus, $W$ esse productum ex $1+i$ in quantitatem realem positivam, unde necessario esse debebit

\[
\begin{aligned}
W=(1+i)\surd n,\quad T=+\surd n,\quad U=+\surd n
\end{aligned}
\]

\Article{19}

Operae pretium erit, omnes summationes hactenus evolutas, hic in unum conspectum colligere.\quad Generaliter scilicet est

\[
\begin{array}{c|c|c}
T=&U=&\text{prout }n\text{ est formae}\\\hline
\pm\surd n&\pm\surd n&4\mu\\
\pm\surd n&0&4\mu+1\\
0&0&4\mu+2\\
0&\pm\surd n&4\mu+3
\end{array}
\]

et in casu eo, ubi $k$ supponitur $=1$, quantitati radicali signum positivum tribui debet.\quad Omni itaque iam rigore ea, quae pro valoribus primis ipsius $n$ in art. 3 per inductionem animadverteramus, demonstrata sunt, nihilque superest, nisi ut signa pro valoribus quibuscunque ipsius $k$ in omnibus casibus determinare doceamus.\quad Sed antequam hoc negotium in omni generalitate aggredi liceat, primo casus eos, ubi $n$ est numerus primus vel numeri primi potestas, propius considerare oportebit.

\Article{20}

Sit primo $n$ numerus primus impar, patetque per ea, quae in art. 10 exposuimus, esse $W=1+2\Sigma r^a=1+2\Sigma R^{ak}$, si statuatur $R=\cos\omega+i\sin\omega$, denotante $a$ ut illic indefinite omnia residua quadratica ipsius $n$ inter $1$ et $n-1$ contenta.\quad Quodsi quoque per $b$ indefinite omnia non-residua quadratica inter eosdem limites exprimimus, nullo negotio perspicitur, omnes numeros $ak$ congruos fieri secundum modulum $n$ vel omnibus $a$ vel omnibus $b$ (nullo ordinis respectu habito), prout $k$ vel residuum sit vel non-residuum.\quad Quamobrem in casu priori erit

\SourceLeaf{42}

\[
\begin{aligned}
W=1+2\Sigma R^a=1+R+R^4+R^9+\text{etc.}+R^{(n-1)^2}
\end{aligned}
\]

adeoque $W=+\surd n$, si $n$ est formae $4\mu+1$, atque $W=+i\surd n$, si $n$ est formae $4\mu+3$.

Contra in casu altero, ubi $k$ est non-residuum ipsius $n$, erit

\[
\begin{aligned}
W=1+2\Sigma R^b
\end{aligned}
\]

Hinc quum manifesto omnes $a$, $b$ complexum integrum numerorum $1$, $2$, $3$ $\ldots$ expleant, adeoque sit

\[
\begin{aligned}
\Sigma R^a+\Sigma R^b=R+R^2+R^3+\text{etc.}+R^{n-1}=-1
\end{aligned}
\]

fiet

\[
\begin{aligned}
W=-1-2\Sigma R^a=-(1+R+R^4+R^9+\text{etc.}+R^{(n-1)^2})
\end{aligned}
\]

adeoque $W=-\surd n$, si $n$ est formae $4\mu+1$, atque $W=-i\surd n$, si $n$ est formae $4\mu+3$.

Hinc itaque colligitur

\emph{primo}, si $n$ est formae $4\mu+1$, atque $k$ residuum quadraticum ipsius $n$,

\[
\begin{aligned}
T=+\surd n,\quad U=0
\end{aligned}
\]

\emph{secundo}, si $n$ est formae $4\mu+1$, atque $k$ non-residuum ipsius $n$,

\[
\begin{aligned}
T=-\surd n,\quad U=0
\end{aligned}
\]

\emph{tertio}, si $n$ est formae $4\mu+3$, atque $k$ residuum ipsius $n$,

\[
\begin{aligned}
T=0,\quad U=+\surd n
\end{aligned}
\]

\emph{quarto}, si $n$ est formae $4\mu+3$, atque $k$ non-residuum ipsius $n$

\[
\begin{aligned}
T=0,\quad U=-\surd n
\end{aligned}
\]
